
\documentclass{article}
\usepackage{colortbl}
\usepackage{makecell}
\usepackage{multirow}
\usepackage{supertabular}

\begin{document}

\newcounter{utterance}

\centering \large Interaction Transcript for game `cladder', experiment `full\_v1.5\_default', episode 6357 with qwen3.5{-}2b.
\vspace{24pt}

{ \footnotesize  \setcounter{utterance}{1}
\setlength{\tabcolsep}{0pt}
\begin{supertabular}{c@{$\;$}|p{.15\linewidth}@{}p{.15\linewidth}p{.15\linewidth}p{.15\linewidth}p{.15\linewidth}p{.15\linewidth}}
    \# & \multicolumn{2}{c}{Player} && \multicolumn{2}{c}{Game Master} \\
    \hline

    \theutterance \stepcounter{utterance}  
    & & & \multicolumn{4}{p{0.6\linewidth}}{
        \cellcolor[rgb]{0.9,0.9,0.9}{
            \makecell[{{p{\linewidth}}}]{
                \texttt{\tiny{[P1$\langle$GM]}}
                \texttt{You are an expert in causal inference. The following question is not a typical commonsense query, but rather a meticulously designed question created by a professor specializing in causal inference, intended to assess the students' mastery of the course content.} \\
\texttt{Imagine a self{-}contained, hypothetical world with only the following conditions, and without any unmentioned factors or causal relationships: Rixq has a direct effect on zuph and xevu. Swoy has a direct effect on zuph and xevu. Zuph has a direct effect on xevu. Swoy is unobserved. The overall probability of rixq is 70\%. The probability of not rixq and xevu is 6\%. The probability of rixq and xevu is 10\%. Is the chance of xevu larger when observing rixq?} \\
\texttt{Start your answer with 'Yes' or 'No', followed by additional reasoning or evidence to support your explanation.} \\
            }
        }
    }
     \\ \\

    \theutterance \stepcounter{utterance}  
    & \multicolumn{4}{p{0.6\linewidth}}{
        \cellcolor[rgb]{0.561,0.737,0.561}{
            \makecell[{{p{\linewidth}}}]{
                \texttt{\tiny{[P1$\rangle$GM]}}
                \texttt{Yes} \\
\\ 
\texttt{**Reasoning:**} \\
\\ 
\texttt{To determine if the chance of \$xevu\$ is larger when observing \$rixq\$, we must calculate the conditional probability \$P(xevu | rixq)\$ and compare it to the marginal probability \$P(xevu)\$.} \\
\\ 
\texttt{1.  **Define the Events and Probabilities:**} \\
\texttt{    *   Let \$R\$ be the event \$rixq\$.} \\
\texttt{    *   Let \$Z\$ be the event \$zuph\$.} \\
\texttt{    *   Let \$E\$ be the event \$xevu\$.} \\
\texttt{    *   Let \$U\$ be the unobserved event \$swoy\$.} \\
\\ 
\texttt{2.  **Analyze the Causal Structure:**} \\
\texttt{    *   \$R \textbackslash{}to Z\$ (Direct effect)} \\
\texttt{    *   \$R \textbackslash{}to E\$ (Direct effect)} \\
\texttt{    *   \$Z \textbackslash{}to E\$ (Direct effect)} \\
\texttt{    *   \$U\$ is unobserved and has no specified effects on \$Z\$ or \$E\$ in the problem description (it is unobserved, implying we cannot use it to adjust for confounding in this specific hypothetical setup, or it is independent).} \\
\\ 
\texttt{3.  **Calculate Marginal Probabilities:**} \\
\texttt{    *   \$P(R) = 0.70\$} \\
\texttt{    *   \$P(\textbackslash{}neg R) = 1 {-} 0.70 = 0.30\$} \\
\texttt{    *   \$P(E | R) = 0.10\$} \\
\texttt{    *   \$P(E | \textbackslash{}neg R) = 0.06\$} \\
\\ 
\texttt{4.  **Calculate \$P(E)\$ (Marginal Probability):**} \\
\texttt{    Using the law of total probability:} \\
\texttt{    \$\$P(E) = P(E | R)P(R) + P(E | \textbackslash{}neg R)P(\textbackslash{}neg R)\$\$} \\
\texttt{    \$\$P(E) = (0.10 \textbackslash{}times 0.70) + (0.06 \textbackslash{}times 0.30)\$\$} \\
\texttt{    \$\$P(E) = 0.07 + 0.018 = 0.088\$\$} \\
\\ 
\texttt{5.  **Calculate \$P(E | R)\$ (Conditional Probability given \$R\$):**} \\
\texttt{    \$\$P(E | R) = 0.10\$\$} \\
\\ 
\texttt{6.  **Compare:**} \\
\texttt{    *   \$P(E | R) = 0.10\$} \\
\texttt{    *   \$P(E) = 0.088\$} \\
\texttt{    *   Since \$0.10 > 0.088\$, observing \$rixq\$ increases the probability of \$xevu\$ compared to the overall baseline probability.} \\
\\ 
\texttt{**Conclusion:**} \\
\texttt{Yes, the chance of \$xevu\$ is larger when observing \$rixq\$ (10\% vs 8.8\%). This is because \$rixq\$ is a direct cause of \$xevu\$, and the probability of \$xevu\$ occurring when \$rixq\$ happens (10\%) is higher than the average probability of \$xevu\$ across all possible states of the world (8.8\%).} \\
            }
        }
    }
    & & \\ \\

    \theutterance \stepcounter{utterance}  
    & & & \multicolumn{2}{p{0.3\linewidth}}{
        \cellcolor[rgb]{0.9,0.9,0.9}{
            \makecell[{{p{\linewidth}}}]{
                \texttt{\tiny{[GM$|$GM]}}
                \texttt{Yes} \\
            }
        }
    }
    & & \\ \\

    \theutterance \stepcounter{utterance}  
    & & & \multicolumn{2}{p{0.3\linewidth}}{
        \cellcolor[rgb]{0.9,0.9,0.9}{
            \makecell[{{p{\linewidth}}}]{
                \texttt{\tiny{[GM$|$GM]}}
                \texttt{no} \\
            }
        }
    }
    & & \\ \\

    \theutterance \stepcounter{utterance}  
    & & & \multicolumn{2}{p{0.3\linewidth}}{
        \cellcolor[rgb]{0.9,0.9,0.9}{
            \makecell[{{p{\linewidth}}}]{
                \texttt{\tiny{[GM$|$GM]}}
                \texttt{game\_result = LOSE} \\
            }
        }
    }
    & & \\ \\

\end{supertabular}
}

\end{document}
